%% kapa: exam=1 round=27 year=2016 period=1 form=A subject=civil_law subject_ko=민법 seq=32 number=32
%% source: Q-Net 감정평가사 기출문제 (cst00309) artlSeq=5169402 file=extracted/1008/5169402/제27회 감정평가사 1차 문제/제27회 감정평가사 1교시문제(A형).pdf
%% conversion: kordoc 4.24.51 -> tools/gen1.py (auto); review=proofread-sonnet; flags=none
\begin{question}{32}
주위토지통행권에 관한 설명으로 옳지 않은 것은? (다툼이 있으면 판례에 따름)
\begin{choices}
  \item 토지의 분할 및 일부양도의 경우, 무상주위통행권에 관한「민법」의 규정은 포위된 토지 또는 피통행지의 특정승계인에게 적용되지 않는다.
  \item 주위토지통행권은 이를 인정할 필요성이 없어지면 당연히 소멸한다.
  \item 기존의 통로가 있더라도 당해 토지의 이용에 부적합하여 실제로 통로로서 충분한 기능을 하지 못하고 있는 경우에도 주위토지통행권이 인정된다.
  \item 통행지소유자는 주위토지통행권자의 허락을 얻어 사실상 통행하고 있는 자에게는 그 손해의 보상을 청구할 수 없다.
  \item 주위토지통행권이 인정되는 도로의 폭과 면적을 정함에 있어서,「건축법」에 건축과 관련하여 도로에 관한 폭 등의 제한규정이 있으면 이에 따라 결정하여야 한다.
\end{choices}
\end{question}
\answer{5}
